\subsection{Algebraic coefficients and a trace-only comparison gate}
VIII allows an algebraic extension of the quantum parameter field.
The single-weight obstruction persists there and requires only
ordinary trace compatibility across extensions.

\begin{finding}[\fid{F51}: the algebraic single-weight trace gate]\label{f:algebraictrace}
Fix $b>1$ and let $\mathcal F$ be an invertible finite matrix
whose coefficients lie in an algebraic extension of $\C(v)$.
Suppose it has a regular invertible value $F_b$ at $v=b$,
and for every sufficiently large integer $n$ choose a regular
value $F_{b^n}$ of the same algebraic family.
Assume
\begin{equation}\label{eq:algebraictraceretune}
 \operatorname{Tr}(F_{b^n})=\operatorname{Tr}(F_b^n).
\end{equation}
If $F_b$ is pure of one fixed integral weight $w$ relative
to $Q=b^2$, then every eigenvalue of $b^{-w}F_b$ is one.
If its full Frobenius centralizer is semisimple, that
centralizer is therefore a single full matrix algebra.
In particular it cannot be the actual generic rank-three
algebra
\[
 A_b\simeq\C\times\C\times M_2(\C)\times M_2(\C).
\]
The conclusion allows different branches of the same
algebraic relation at different extension parameters;
it assumes neither rational coefficients nor full
spectral retuning.

More generally, suppose the base eigenvalues have several
integral weights $w_j$, with genuine multiplicities.
The same trace comparison forces the family trace to be
the corresponding positive Laurent polynomial and every
base eigenvalue to be $b^{w_j}$, without a nontrivial phase.
If full spectral retuning is imposed as well, then
\[
 \det(TI-\mathcal F)=\prod_w(T-v^w)^{m_w}.
\]
A regular fixed-rank model at $v=1$ therefore cannot
have the actual collision algebra as its full special
Frobenius centralizer.
\end{finding}

\begin{proof}
We first prove the precise algebraic-function statement.
Let
\[
 u_n=\sum_{j=1}^s c_j\alpha_j^n,\qquad |\alpha_j|=1,
\]
with finitely many roots and arbitrary complex coefficients.
Let $P(v,Y)\in\C[v,Y]$ be irreducible of positive
$Y$-degree, and suppose
\[
 P(b^n,u_n)=0
\]
for every sufficiently large integer $n$.
We claim that
\begin{equation}\label{eq:radialconstant}
 P(v,Y)=a(Y-u_0),\qquad a\ne0,
\end{equation}
and $u_n=u_0$ for every integer $n$.

Write
\[
 P(v,Y)=\sum_{i=0}^d v^iP_i(Y).
\]
Each polynomial $P_i(u_n)$ is a finite power sum on
unit-circle roots: after expanding its powers and grouping
equal products,
\[
 P_i(u_n)=\sum_{\beta\in B_i} a_{i,\beta}\beta^n,
 \qquad |\beta|=1.
\]
Its constant term is represented by the root $\beta=1$.
The assumed relation becomes
\[
 \sum_i\sum_{\beta\in B_i}
       a_{i,\beta}(b^i\beta)^n=0.
\]
Different $i$ give different root moduli $b^i$, so no
roots from those groups coincide.
Within a group equal roots have already been collected.
Evaluating a relation on distinct nonzero power roots
at sufficiently many consecutive integers gives an
invertible Vandermonde matrix times an invertible
diagonal matrix. Thus every $a_{i,\beta}=0$.

Consequently every $P_i(u_n)$ vanishes as an exponential
sequence for all integers, including $n=0$.
Therefore $P_i(u_0)=0$ for every $i$, or
$P(v,u_0)=0$ identically in $v$.
The factor $Y-u_0$ divides $P$ in $\C[v,Y]$.
Irreducibility proves \eqref{eq:radialconstant}.
The eventual equality $u_n=u_0$ then extends to all
integers by the same finite-power-sum uniqueness.
The auxiliary value $n=0$ is an algebraic sequence value,
not an assertion of an extension of degree zero or a
field of cardinality one.

For a genuine $N$-dimensional matrix spectrum,
$u_n=\operatorname{Tr}(G^n)$ and $u_0=N$.
If all eigenvalues of $G$ have modulus one, the preceding
argument gives $u_n=N$ eventually.
A sum of $N$ unit-modulus numbers equals the real number
$N$ only when each summand is one.
Two consecutive sufficiently large $n$ then imply
$\alpha_j^n=\alpha_j^{n+1}=1$, hence $\alpha_j=1$
for every eigenvalue.
This is where genuine positive multiplicities are used.
For a virtual trace the same argument concerns only
its visible grouped roots.

Return to the family and normalize
\[
 \mathcal G=v^{-w}\mathcal F,\qquad G_b=b^{-w}F_b.
\]
The trace $g=\operatorname{Tr}\mathcal G$ is algebraic
over $\C(v)$. Its irreducible minimal polynomial,
after clearing denominators and taking its primitive
part, is an irreducible $P(v,Y)\in\C[v,Y]$.
Every regular evaluation of this same algebraic element
satisfies that relation.
By \eqref{eq:algebraictraceretune},
\[
 g_{b^n}
 =b^{-nw}\operatorname{Tr}(F_b^n)
 =\operatorname{Tr}(G_b^n).
\]
The eigenvalues of $G_b$ all have modulus one.
The preceding argument proves $g(v)=N$ and that every
eigenvalue of $G_b$ is one. Thus $G_b-I$ is nilpotent.

Scalar normalization does not change the full centralizer.
The nilpotent $G_b-I$ is central in that algebra.
If the algebra is semisimple, its center is reduced,
so $G_b-I=0$.
The full centralizer is then $M_N(\C)$.
The actual algebra of \fid{F19} at $b>1$ has its two
center branches split and is the nonsimple semisimple
product displayed in the statement. It cannot be
such a full centralizer.

\par\smallskip\noindent\textbf{Several integral spectral weights.}\quad
We prove the more general algebraic-function assertion
with arbitrary finite complex coefficients first.
Write
\[
 u_n=\sum_j c_j b^{w_jn}\zeta_j^n,\qquad
 w_j\in\Z,\quad |\zeta_j|=1,
 \qquad p(v)=\sum_j c_jv^{w_j}.
\]
Suppose $P(b^n,u_n)=0$ eventually for the irreducible
minimal relation of one algebraic function $f$.
Expand
\[
 P(v,Y)=\sum_{i,k}a_{i,k}v^iY^k.
\]
Every root in its evaluated power-sum expansion has form
\[
 b^{\,i+w_{j_1}+\cdots+w_{j_k}}
       \zeta_{j_1}\cdots\zeta_{j_k}.
\]
Group these by the integer
$r=i+w_{j_1}+\cdots+w_{j_k}$.
For fixed $r$ the coefficient is a finite unit-circle
power sum $S_r(n)$.
Distinct $r$ give distinct root moduli $b^r$, so
Vandermonde uniqueness makes each $S_r(n)$ identically
zero. At $n=0$ all phases disappear.
The value $S_r(0)$ is exactly the coefficient of $v^r$
in $P(v,p(v))$. Hence
\[
 P(v,p(v))=0\quad\hbox{in }\C[v,v^{-1}].
\]
Irreducibility over $\C(v)$ and the rational root $p(v)$
force $f=p$.
This argument is taken over $\C(v)$ when weights are
negative; no polynomial factorization with an omitted
denominator is being assumed.

Thus $u_n=p(b^n)$ eventually and, by finite-power
uniqueness, identically as a sequence.
For genuine positive eigenvalue multiplicities, any
nontrivial phase would give a positive coefficient on
a root absent from the Laurent side. It cannot cancel.
Every base eigenvalue is therefore $b^w$, with its
stated weight multiplicity.

\par\smallskip\noindent\textbf{Full spectra and the collision.}\quad
Now additionally assume full spectral retuning.
Let $a(v)$ be an algebraic characteristic-polynomial
coefficient, with irreducible minimal relation $P_a$.
Its evaluated values agree with the corresponding
Laurent coefficient $p_a(v)$ of
$\prod_w(T-v^w)^{m_w}$.
Thus $P_a(b^n,p_a(b^n))=0$ at infinitely many points.
The rational function $P_a(v,p_a(v))$ is identically
zero, and its rational root forces $a=p_a$.
This proves the stated characteristic-polynomial identity,
even when different allowed sheets are chosen at
different extension parameters.

If the family is regular at a point above $v=1$,
its characteristic polynomial there is $(T-1)^N$.
The center of its full centralizer is therefore local,
by the complete cyclic-module proof in
\eqref{eq:nilpotentcentralizercenter}.
The actual center
$k\times k\times k[\epsilon]/(\epsilon^2)$
has three components, giving the collision contradiction.
\end{proof}

\paragraph{Algebraic sheets and the exact scope.}
The finitely many matrix coefficients generate a finite
extension even if the initial coefficient extension is
not finite. Their trace is one algebraic element.
Allowed evaluations may choose different sheets at
different $b^n$, but must satisfy its same irreducible
relation. Switching to an unrelated algebraic function
or a different constant-field component is not part
of that assumption.
No continuation at infinity, common eigenvalue labeling,
or diagonalization of the family was used.
The trace and single-weight conclusions require no
regularity at $v=1$.
For the inverse quantum-parameter convention, replace
$v$ by $v^{-1}$; the normalization becomes $v^w$
in the original coordinate.

From trace comparison alone, only the base operator
is proved unipotent after single-weight normalization.
Determining the full characteristic polynomial at
other parameters above used the separately stated
full spectral comparison.
The obstruction permits any auxiliary rank and Jordan
blocks before the comparison is imposed.
It addresses one fixed integral pure weight and a
finite genuine spectrum, with the ordinary trace
rule \eqref{eq:algebraictraceretune}.
The several-weight extension requires integral spectral
weights. The collision conclusion additionally requires
full spectral retuning and a regular fixed-rank model.
This extends the rational collision theorem of
\fid{F48} to algebraic entries in that integral-weight
case. Nonintegral spectral weights are not classified
by this argument.
Infinite, non-atomic, regularized and virtual traces
require their own convergence and visibility checks.

The fixed-parameter models of \fid{F46} and the
faithful tensor actions of \fid{F48} remain valid.
What is incompatible is their proposed strengthening
to a full nonsimple semisimple centralizer, one pure
weight, an algebraic parameter family and this
all-extension trace comparison.
This directly covers the algebraic coefficient-field
format allowed in VIII, \S2.2, p.~7 \cite{VIII},
when that additional comparison is required.
It is not a general nonstandard Riemann hypothesis.

% Separate paragraph for Part III.
\paragraph{\fid{F51} in Grothendieck's language.}
The coefficient extension is the algebraic field
extension explicitly allowed in VIII, \S2.2.
The module and action interpretation uses EGA~I,
Theorem~1.3.7, p.~86, Corollary~1.3.8, p.~87, and
Theorem~1.4.1, pp.~90--92; the coaction framework is
SGA~3, I, 4.7.2--3, with cote~90, scan~17 as the
audited anchor.
The pure-weight convention is BBD~5.1.2 and~5.1.8.
These frameworks do not assert the arithmetic
comparison: the displayed polynomial, radius and
Vandermonde argument proves its particular restriction.
No unpublished archive leaf is claimed to contain
this algebraic trace obstruction.

